% !TEX root = ../sparsity_thesis.tex
\chapter{Conclusion}
\label{ch:conclusion}

\section{Conclusion}
\label{sec:conclusion}

This research set out to investigate a structural blind spot in the
deployment of Explainable Artificial Intelligence within multi-tenant
Enterprise Resource Planning systems: the \emph{hallucinated explanation}
phenomenon, in which post-hoc explainers continue to produce confidently
formatted, visually emphatic outputs even when applied to models trained on
data-sparse tenants whose underlying predictive signal has collapsed. The
empirical evidence gathered in this study demonstrates that data sparsity
degrades explanation fidelity in a non-linear and qualitatively misleading
manner. Across the five simulated tenant tiers, and under a strict
leakage-avoidance regime in which only order-time observables were admitted
into the feature matrix, SHAP rank correlation against the data-rich
baseline fell monotonically from $1.0000$ to $0.4485$ and crossed below the
conventional reliability bound of $0.70$ at the $n = 1{,}000$ tier, while
SHAP absolute variance simultaneously rose by more than two orders of
magnitude from $0.0224$ at the data-rich tier to $5.0392$ at the cold-start
tier. By contrast, classical predictive metrics did not reveal this regime
change: AUC-ROC remained uniformly modest across tiers~A--D ($\approx 0.61$
--$0.66$) and anomalously \emph{rose} to $0.8889$ at the cold-start tier
as a small-sample artefact. Taken
together, these findings establish that XAI applied to cold-start tenants
is not merely \emph{less accurate}---it is \emph{unfaithful} in rank order
and \emph{volatile} in attribution magnitude, while remaining
indistinguishable from a reliable explanation in its rendered visual form.
This directly answers Research Question~1.

The study further establishes that the operational reliability of local
post-hoc explanations exhibits a regime change---an empirical
\emph{Reliability Cliff}---at the $n = 1{,}000$-record threshold. This
boundary is not the product of a single arbitrary metric; rather, it is
supported by the simultaneous convergence of two independent
explanation-fidelity indicators at the same tier: the Spearman
rank-correlation criterion ($\rho_C = 0.6429 < 0.70$) and the SHAP
variance criterion ($\sigma^2_{|\mathrm{SHAP}|,C}$ rising $128\%$ above
the data-rich baseline). The classical AUC criterion does not provide
additional support in the leakage-free regime---predictive performance is
uniformly weak across tiers~A--D and anomalously strong at the cold-start
tier---which makes the explanation-fidelity criteria \emph{not merely
convergent with, but operationally necessary for}, the localisation of the
cliff. Below this threshold, the underlying machine-learning models and
their post-hoc explanations diverge from any genuine business logic and
cannot be trusted for operational decision-making. The cross-task
replication on the Employee Attrition and Lead Scoring tasks
(Section~\ref{sec:cross-task-replication}) establishes that the
qualitative regime change is not an artefact of the Late Delivery Risk substrate:
the SHAP-variance expansion and Spearman degradation recur on both
additional ERP-native tasks, while the absolute threshold $\tau$ at
which the cliff occurs is task-specific. This directly answers
Research Question~2.

The study concludes that the conventional practice of deploying a uniform
machine-learning and XAI pipeline across every tenant of a multi-tenant
SaaS platform is fundamentally flawed, because it ignores the per-tenant
data-maturity distribution that is intrinsic to the SaaS business model.
The \emph{Maturity-Aware AI Routing} pattern proposed and implemented in
this research provides the necessary architectural correction. By placing
the \texttt{predict\_or\_heuristic()} gate inside the
\texttt{LatePaymentPredictor} service of the studied Django ERP, and by
parameterising it on the empirically-derived
\texttt{COLD\_START\_INVOICE\_THRESHOLD = 1000}, the system structurally
protects cold-start tenants from being shown hallucinated local
explanations, serving them platform-wide Global Heuristics drawn from
data-rich SHAP analyses instead, while continuing to deliver granular
local SHAP and LIME explanations to data-mature tenants for whom the
empirical evidence supports their reliability. This directly answers
Research Question~3.

The principal academic takeaway of this thesis is therefore that
Explainable Artificial Intelligence must not be conceived of as a passive
reporting feature appended to a model after the fact, but as an active
\emph{diagnostic instrument} that exposes the limits of model reliability
and conditions the way intelligence is delivered to the end user. In
multi-tenant SaaS platforms, where data maturity is heterogeneous by
construction, this reframing is not optional: it fundamentally shifts the
responsibility of the system architect from \emph{displaying} explanations
to \emph{governing} them, and it establishes data maturity as a
first-class routing dimension in any production AI subsystem.

\section{Suggestions}
\label{sec:suggestions}

The findings of this study open several productive directions for further
work, which may be organised into suggestions for future academic research
and suggestions for further system development.

\subsection{For Future Research (Academic)}
\label{subsec:future-research}

\begin{enumerate}
    \item Future studies should investigate XAI fidelity degradation in
    \emph{longitudinal} and \emph{live} SaaS environments, rather than the
    cross-sectional simulated datasets used in this thesis. A longitudinal
    design would allow the interaction between data sparsity and
    \emph{concept drift} to be characterised, addressing the realistic case
    in which a tenant's data distribution itself evolves as the tenant
    matures.
    \item The cold-start XAI problem should be extended beyond the tabular
    gradient-boosting setting studied here to \emph{deep learning} model
    families, in particular recurrent and attention-based architectures
    (e.g., LSTMs and Transformers) applied to time-series ERP forecasting,
    together with their attendant explainers (Integrated Gradients,
    Attention rollout). The quantitative threshold derived in this thesis
    is unlikely to transfer literally to these settings and would need to
    be re-established empirically.
    \item Controlled user studies should be conducted with business
    operators of multi-tenant ERPs in order to measure the \emph{perceived
    trust} and \emph{decision-making accuracy} induced by Global Heuristics
    versus Local Explanations in cold-start scenarios. This would
    complement the present technical evidence of explanation unfaithfulness
    with behavioural evidence of its downstream consequences on end-user
    judgement.
\end{enumerate}

\subsection{For System Development (Practical)}
\label{subsec:system-development}

\begin{enumerate}
    \item The static \texttt{COLD\_START\_INVOICE\_THRESHOLD = 1000}
    constant should be replaced, in production deployments, by a
    \emph{dynamic threshold mechanism} that auto-adjusts in response to
    real-time data-distribution indicators---for example,
    Kolmogorov--Smirnov drift statistics computed against the global
    reference distribution, or per-tenant data-velocity
    (records-per-week) metrics. This would allow the routing gate to
    remain calibrated as both the platform and the individual tenant
    evolve.
    \item The Maturity-Aware AI Routing pattern should be generalised
    beyond the Late Delivery Risk task to the other AI-supported modules of the
    ERP (Lead Scoring, Employee Attrition, Procurement Risk,
    Expense-Claim Anomaly). Each module is expected to exhibit a distinct
    empirical threshold reflecting its task complexity and class
    prevalence, and the routing layer should accept per-module
    configuration accordingly.
    \item An interactive \emph{Data Maturity Dashboard} should be
    developed for ERP administrators, visually communicating each
    tenant's progress toward the data-maturity threshold required for
    personalised local AI predictions. Such a dashboard converts the
    architectural threshold from an invisible engineering constant into a
    transparent business-facing artefact, supports onboarding planning,
    and aligns end-user expectations with the technical preconditions of
    trustworthy local explanations.
\end{enumerate}
